\documentclass[12pt,a4paper]{article}
\usepackage{fontspec}
\usepackage{polyglossia}
\setmainfont{TeX Gyre Termes}
\newfontfamily\cyrillicfont{TeX Gyre Termes}
\setdefaultlanguage{ukrainian}
\setotherlanguage{english}
\usepackage{amsmath,amssymb,geometry,graphicx}
\geometry{margin=25mm}
\setlength{\parindent}{1.25cm}
\setlength{\parskip}{0.25em}
\begin{document}
\title{Tezy 2025 Переваги методу політочкових перетворень порівняно з metодами NeRF та gaussian splatting}
\maketitle
ПЕРЕВАГИ МЕТОДУ ПОЛІТОЧКОВИХ ПЕРЕТВОРЕНЬ ПОРІВНЯНО З METОДАМИ NeRF ТА GAUSSIAN SPLATTING Актуальність досліджень у сфері комп’ютерної графіки та геометричного моделювання зумовлена широким застосуванням технологій деформації об’єктів у системах віртуальної реальності, кінематографі, CAD/CAM системах та навчальних симуляторах. Однак сучасні методи, такі як Neur al Radiance Fields (NeRF) та Gaussian Splatting, хоча й забезпечують високу якість візуалізації, мають суттєві обмеження, особливо щодо точності керування формою та реального часу застосування. Визначено основні недоліки методу NeRF, серед яких виділяється складність керування геометричними змінами через необхідність великого обсягу обчислювальних ресурсів і навчання нейронної мережі для кожного конкретного випадку. Метод Gaussian Splatting, хоча й дозволяє швидку та ефективну реалізацію у реальному часі, х арактеризується зниженою точністю локальних деформацій, оскільки використовує дискретні гауссові ядра, що ускладнює точну локалізацію змін форми об'єкта. На противагу цим методам, метод політочкових перетворень забезпечує низку вагомих переваг. Його базисо м є скінченна множина точок, що дозволяє ефективно керувати локальними деформаціями без значних обчислювальних витрат. Математичний апарат цього методу використовує систему лінійних рівнянь, що гарантує простоту та прозорість розрахунків, на відміну від «чорного ящика» нейронних методів. Також політочкові перетворення дозволяють уникати проблеми взаємної неоднозначності, характерної для багатьох інших методів, шляхом застосування спеціально сформульованих додаткових умов мінімального відхилення. Таким чином, метод політочкових перетворень є перспективною альтернативою сучасним нейронним методам деформації, оскільки пропонує інтуїтивно зрозумілий інструментарій для керування формою об’єктів, забезпечує високу точність локальних змін і має потенціал для ефективної реалізації у реальному часі, що робить його придатним для широкого спектра прикладних задач у сфері комп’ютерної графіки та моделювання.
\par
\end{document}
